\documentclass[11pt, a4paper]{article}

% --- Packages ---
\usepackage[utf8]{inputenc}
\usepackage[T1]{fontenc}
\usepackage[french]{babel}
\usepackage[left=1.5cm, right=1.5cm, top=1.5cm, bottom=1.5cm]{geometry}
% \usepackage{fontawesome5} % Pour les icônes (temporarily disabled)
\usepackage{hyperref}     % Pour les liens cliquables
\usepackage{titlesec}     % Pour personnaliser les titres de section
\usepackage{enumitem}     % Pour personnaliser les listes
\usepackage{parskip}      % Pour espacer les paragraphes
\usepackage{xcolor}

% --- Configuration des couleurs et liens ---
\definecolor{linkcolor}{rgb}{0.0, 0.0, 0.5} % Bleu sombre et pro pour les liens
\hypersetup{
    colorlinks=true,
    linkcolor=linkcolor,
    urlcolor=linkcolor
}

% --- Configuration des titres de sections ---
\titleformat{\section}{\large\scshape\raggedright\bfseries}{}{0em}{}[\titlerule]
\titlespacing{\section}{0pt}{10pt}{5pt}

% --- Commandes personnalisées pour le CV ---
% Commande pour une entrée classique (Expérience / Formation)
\newcommand{\cventry}[4]{
    \vspace{0.15cm}
    \noindent \textbf{#1} \hfill #2 \\
    \noindent \textit{#3} \hfill \textit{#4} \par
    \vspace{0.1cm}
}

\begin{document}
\pagestyle{empty} % Pas de numéro de page

% ================= EN-TÊTE =================
\begin{center}
    {\Huge \textbf{Marcel ASSIE}} \\[0.2cm]
    {\Large \textsc{Ingénierie Géomatique}} \\[0.3cm]
    \small
    Tel: 07 68 06 25 78 \quad|\quad
    Email: \href{mailto:marcelassie2k@gmail.com}{marcelassie2k@gmail.com} \quad|\quad
    LinkedIn: \href{https://www.linkedin.com/in/marcel-assie/}{LinkedIn} \quad|\quad
    GitHub: \href{https://github.com/MarcelAssie/}{GitHub} \\[0.15cm]
    Web: \href{https://marcel-assie-portefolio.dev}{marcel-assie-portefolio.dev} \quad|\quad
    Location: Champs-sur-Marne (77), France
\end{center}

\vspace{0.3cm}
\noindent À la recherche d’une opportunité en \textbf{CDI} pour mettre à profit mes compétences en développement géospatial, analyse de données spatiales et automatisation, \textbf{au service de projets data} à fort impact.

% ================= FORMATIONS =================
\section*{Formations}

\cventry{École Nationale des Sciences Géographiques (Geodata Paris)}{Champs-sur-Marne, France}{Ingénierie Géomatique}{2023 -- 2026}
\cventry{Institut National Polytechnique Félix Houphouët-Boigny}{Yamoussoukro, Côte d'Ivoire}{Licence Géomètre}{2019 -- 2022}
\cventry{Lycée Scientifique de Yamoussoukro}{Yamoussoukro, Côte d'Ivoire}{Baccalauréat, Série Scientifique}{2016 -- 2019}

% ================= EXPÉRIENCES PROFESSIONNELLES =================
\section*{Expériences Professionnelles}

\cventry{ISOGEO}{Paris, France}{Ingénieur Développeur SIG -- Alternance}{Depuis Sept. 2025}
\begin{itemize}[leftmargin=0.5cm, nosep]
    \item Remplacement et optimisation de flux FME par des solutions alternatives (Python, SQL).
    \item Développement de scripts d’automatisation pour l’extraction et la signature de métadonnées (Python).
    \item Étude et intégration de modèles d’IA pour l’enrichissement automatique de métadonnées.
    \item Rédaction de spécifications techniques et contribution aux développements en méthodologie Agile.
\end{itemize}

\cventry{ORYJIN}{Lille, France}{Ingénieur Geo Data Scientist -- Stage}{Avril 2025 -- Août 2025}
\begin{itemize}[leftmargin=0.5cm, nosep]
    \item Développement de pipelines ETL (Extract, Transform, Load) pour la préparation et l’agrégation de données (Python).
    \item Automatisation de l’intégration des données dans un référentiel spatial basé sur H3 (Python, Snowflake, SQL).
    \item Mise en place d’algorithmes d’analyse de photographies aériennes à l’aide de modèles d’IA (Python).
    \item Prompt engineering pour la génération de visuels d’aide à la décision.
\end{itemize}

\cventry{IGN}{Forcalquier, France}{Ingénieur Géomaticien -- Stage}{Mai 2024 -- Juillet 2024}
\begin{itemize}[leftmargin=0.5cm, nosep]
    \item Mini projets appliqués en topométrie, photogrammétrie, métrologie (auscultation de pont), télédétection et géodésie, avec formalisation méthodologique et analyse des incertitudes.
    \item Conception d’un workflow automatisé de cartographie des restrictions de circulation (Alpes-de-Haute-Provence) :
    \begin{itemize}[leftmargin=0.5cm, nosep]
        \item Structuration et nettoyage de jeux de données multi-sources.
        \item Modélisation des traitements géospatiaux (QGIS Model Builder).
        \item Développement d’algorithmes de contrôle qualité et de validation des données.
    \end{itemize}
\end{itemize}

% ================= COMPÉTENCES =================
\section*{Compétences}
\noindent \textbf{Informatique / Programmation :} Python, JavaScript, PHP, SQL \\
\textbf{Technologies \& Outils :} Docker, GitHub, PostgreSQL, PostGIS, MySQL, Django, Flask, Vue.js, React Native, Neo4j \\
\textbf{Outils SIG :} QGIS, ArcGIS, FME, GeoServer, Leaflet \\
\textbf{Aptitudes professionnelles :} Forte capacité d’adaptation, Rigueur, Esprit critique et d'analyse, Prise d’initiative

% ================= PROJETS =================
\section*{Projets}

\noindent \textbf{Projets Académiques}
\begin{itemize}[leftmargin=0.5cm, nosep]
    \item \textbf{Extraction de données contextuelles (ODD)} : Utilisation de LLM et graphes de connaissances (Python, NLP, Jupyter, Neo4j, Cypher, modèles d’IA).
    \item \textbf{App mobile de navigation (Gendarmerie Nationale)} : Application de type Waze (React Native, SQL, JS, Flask).
    \item \textbf{Étude Spatio-Statistique (Île-de-France)} : Analyse des flux domicile-travail et déterminants du choix modal (Statistiques, Machine Learning, Analyse spatiale).
\end{itemize}

\vspace{0.2cm}
\noindent \textbf{Projets Personnels}
\begin{itemize}[leftmargin=0.5cm, nosep]
    \item \textbf{Plateforme SaaS KlassIvoire} : Contribution au développement pour le suivi scolaire numérique, avec intégration d’outils d'IA pour l'analyse des performances (Django, Vue.js, Leaflet, SQL).
    \item \textbf{Plateforme d'analyse d'établissements scolaires} : Conception d’une plateforme intelligente de suivi géographique permettant l’exploitation de données spatiales pour la prise de décision.
\end{itemize}

% ================= VIE ASSOCIATIVE =================
\section*{Vie Associative}

\cventry{VERTIGEO -- Junior Entreprise de l'ENSG Géomatique}{Champs-sur-Marne}{Secrétaire Général}{Fév. 2024 -- Mars 2025}
\noindent Comptes rendus, archivage, gestion des relations avec les partenaires administratifs, organisation de réunions.

\cventry{Professeur particulier}{}{Cours à domicile -- Mathématiques et Physique}{Depuis Oct. 2023}
\noindent Pédagogie et méthodologie, écoute attentive des besoins des élèves.

\cventry{Épicerie solidaire campésienne}{}{Bénévole}{Fév. 2024}
\noindent Collecte de produits alimentaires et d'hygiène à Carrefour. Promotion de l'échange et de l'ouverture d'esprit.

% ================= LANGUES ET INTÉRÊTS =================
\section*{Langues et Centres d'intérêt}
\noindent \textbf{Langues :} Français (Langue maternelle), Anglais (TOEIC - B2) \\[0.1cm]
\textbf{Centres d'intérêt :} Intelligence Artificielle, Veille technologique, Football, Gymnastique acrobatique, Jeux de société

\end{document}